% ============================================================================
% 09_modelo_datos.tex — Cuerpo del entregable "Modelo de Base de Datos"
% Sin preámbulo: se incluye tanto desde 09_Modelo_Datos.tex (standalone)
% como, eventualmente, desde un documento maestro.
% ============================================================================
\section{Modelo de Base de Datos}

\subsection{Justificación del modelo de persistencia híbrido}

LSA Salud adopta una \textbf{arquitectura de persistencia híbrida}, que combina
un motor de base de datos relacional embebido (\textbf{SQLite}) para la
información estructurada del sistema con el \textbf{sistema de archivos local}
para los artefactos binarios de aprendizaje automático y procesamiento de audio.

Esta decisión de diseño responde a las siguientes razones técnicas y
arquitectónicas, verificables contra el código fuente y contra los requisitos
de la ERS:

\begin{itemize}
    \item \textbf{Motor relacional embebido (SQLite en \texttt{data/lsa\_salud.db}).}
    Los reportes de traducción (RF-06.1) y los eventos de auditoría y métricas
    (RF-04.1) son registros estructurados que exigen integridad referencial,
    consistencia transaccional (ACID) y consultas analíticas agregadas (por
    ejemplo, distribución de reportes por palabra, latencias percentiladas y
    conteo de eventos por estación). SQLite resuelve estos requerimientos de
    forma nativa en un único archivo local mediante transacciones protegidas
    por WAL (\textit{Write-Ahead Logging}), garantizando durabilidad y
    consultas indexadas sin requerir servicios externos ni demonios
    adicionales.

    \item \textbf{Localidad y autonomía (RNF-T-02 y RNF-03.1).} Cada estación
    fija opera de forma autónoma en una PC estándar sin conexión a servidores
    centrales ni puertos de base de datos abiertos hacia la red. Un servidor de
    base de datos relacional externo (como PostgreSQL o MySQL) introduciría una
    superficie de ataque innecesaria, dependencias de infraestructura y
    consumo de memoria ocioso en el puesto de atención. SQLite reside en el
    mismo proceso del backend Python (\texttt{import sqlite3}), garantizando
    que ningún dato abandone la máquina física.

    \item \textbf{Formato nativo para tensores de aprendizaje automático.} Las
    muestras de señas se capturan y consumen como arreglos multidimensionales
    de NumPy (\texttt{float32}) de forma $(45, 75, 4)$ para landmarks y
    $(45, 52)$ para blendshapes faciales; los pesos del modelo corresponden a un
    \texttt{state\_dict} de PyTorch y la voz sintética a un grafo ONNX.
    Almacenar estos blobs binarios en tablas relacionales solo añadiría capas
    de serialización y sobrecarga de I/O sin aportar capacidades de consulta.
    Por ello, se mantienen directamente en el sistema de archivos
    (\texttt{dataset/}, \texttt{models/}, \texttt{voices/}).

    \item \textbf{Política estricta sobre la conversación médica (RNF-03.1).}
    El historial de conversación ordinario entre el paciente y el personal de
    recepción o enfermería \textbf{no se persiste en disco}. Vive
    exclusivamente en la memoria del navegador como estado de React sincronizado
    en tiempo real entre pestañas mediante \texttt{BroadcastChannel("nexora-chat")}
    (ver \texttt{web/hooks/useChatChannel.ts}). Al cerrar la pestaña o al pulsar
    el botón de finalización de atención (CU-12), la memoria se libera por
    completo. Esto evita acumular información médica sensible innecesaria.
\end{itemize}

% ----------------------------------------------------------------------------
\subsection{¿Dónde y cómo persiste el mensaje?}
\label{sec:dat-donde-persiste}

Uno de los puntos críticos del diseño de datos de LSA Salud es la respuesta
precisa a la persistencia del mensaje:

\begin{enumerate}[leftmargin=*]
    \item \textbf{En el flujo normal de atención (CU-01 y CU-04):} el mensaje
    \textbf{no se persiste en ninguna base de datos ni archivo}. El texto
    traducido o dictado se envía mediante el canal \texttt{BroadcastChannel}
    de la API web hacia las pestañas de ambas caras de la estación. Reside
    en memoria RAM del cliente y desaparece cuando concluye la sesión.
    
    \item \textbf{En el flujo de reporte de error o ambigüedad (CU-10):} aquí se
    produce la \textbf{única excepción controlada} a la regla de no
    persistencia. Cuando el paciente o el personal hacen clic en ``Reportar
    traducción'', el sistema envía una petición \texttt{POST /report} al backend.
    En ese instante exacto:
    \begin{itemize}[leftmargin=*, itemsep=2pt]
        \item El texto del mensaje señalado se inserta en la tabla relacional
        \textbf{\texttt{reported\_message}} dentro de \texttt{data/lsa\_salud.db}.
        Si ambas caras de la estación reportan el mismo mensaje, la restricción
        de unicidad \texttt{UNIQUE (client\_message\_id, text)} evita
        duplicar el contenido textual, vinculando ambos reportes a la misma fila.
        \item Las palabras reconocidas de la seña que dieron origen a la
        traducción se registran en la tabla secundaria
        \textbf{\texttt{recognized\_word}}, asociadas por clave foránea al
        mensaje (\texttt{ON DELETE CASCADE}).
        \item Cada acción de reporte individual se graba en la tabla
        \textbf{\texttt{translation\_report}}, indicando qué estación
        (\texttt{paciente} o \texttt{recepcion}) emitió el reclamo y la marca
        temporal exacta.
        \item En el frontend, el mensaje se marca como reportado en ambas caras
        de la estación, se exhibe tachado (\textit{strikethrough}) y se
        solicita al paciente la repetición de la seña, evitando que el
        personal tome como válida una traducción equívoca.
    \end{itemize}
\end{enumerate}

De este modo, LSA Salud cumple de forma simultánea con el principio de
minimización de datos clínicos (RNF-03.1) y con la necesidad de ingeniería de
capturar los casos de falla reales como insumo de reentrenamiento del modelo
de IA (RF-06.1 y CU-08).

% ----------------------------------------------------------------------------
\subsection{Modelo conceptual (Entidad-Relación)}

La Figura~\ref{fig:dat-er-conceptual} presenta el modelo conceptual de datos de
LSA Salud. Se aprecian con claridad tres subsistemas: los artefactos binarios
del sistema de archivos, las tablas de la base de datos relacional SQLite y la
representación volátil del mensaje en el navegador.

\diagrama{09_er_conceptual}{\textwidth}{Diagrama Entidad-Relación conceptual de la persistencia de LSA Salud}{fig:dat-er-conceptual}

El modelo explicita las siguientes relaciones esenciales:
\begin{itemize}[leftmargin=*, itemsep=2pt]
    \item \textbf{\textit{Mensaje} $\rightarrow$ \textit{MensajeReportado} (0..1):}
    representa el paso de la memoria volátil a disco. De todos los mensajes
    emitidos en una consulta médica, solo aquellos que son objeto de un reporte
    llegan a grabarse en la base de datos.
    \item \textbf{\textit{MensajeReportado} $\rightarrow$ \textit{ReporteDeTraducción} (1..*):}
    un mensaje reportado cuenta con al menos un reporte y puede acumular
    múltiples reportes si es señalado de forma independiente por el paciente y
    por el personal de recepción.
    \item \textbf{\textit{MensajeReportado} $\rightarrow$ \textit{PalabraReconocida} (0..20):}
    composición que registra los términos crudos identificados por el clasificador
    antes de la síntesis gramatical.
    \item \textbf{\textit{PalabraReconocida} $\rightarrow$ \textit{ClaseDeSeña} (0..1):}
    correspondencia semántica con las clases del vocabulario de entrenamiento
    que permite al equipo de IA priorizar qué señas requieren más muestras.
    \item \textbf{\textit{EventoDeAuditoría}:} entidad independiente para
    trazabilidad y métricas de latencia y consentimiento (RF-03/04).
\end{itemize}

% ----------------------------------------------------------------------------
\subsection{Modelo lógico}

La Figura~\ref{fig:dat-modelo-logico} traduce el modelo conceptual a un esquema
lógico con claves primarias (\texttt{PK}), foráneas (\texttt{FK}), tipos de datos,
restricciones y cardinalidades.

\diagrama{09_modelo_logico}{\textwidth}{Modelo lógico de la persistencia de LSA Salud}{fig:dat-modelo-logico}

Las tablas de SQLite implementan las siguientes restricciones relacionales:
\begin{itemize}[leftmargin=*, itemsep=2pt]
    \item \textbf{\texttt{reported\_message}:} clave primaria autoincremental
    \texttt{id}. Restricción \texttt{UNIQUE (client\_message\_id, text)} para
    evitar duplicaciones al sincronizar entre las dos caras de la estación.
    \item \textbf{\texttt{recognized\_word}:} clave primaria compuesta
    (\texttt{message\_id}, \texttt{position}) que garantiza el orden exacto de
    las palabras, con clave foránea referenciando a \texttt{reported\_message(id)}
    con política \texttt{ON DELETE CASCADE}.
    \item \textbf{\texttt{translation\_report}:} clave foránea indexada
    \texttt{message\_id} y validación \texttt{CHECK (station IN ('paciente', 'recepcion'))}.
    \item \textbf{\texttt{audit\_event}:} tabla con índice compuesto
    \texttt{(type, occurred\_at)} para resolver de forma instantánea las
    consultas analíticas por rango de fechas de \texttt{GET /metrics}.
\end{itemize}

% ----------------------------------------------------------------------------
\subsection{Diccionario de datos}

Las Tablas~\ref{tab:dat-dic-reported-msg} a~\ref{tab:dat-dic-conversacion}
detallan los atributos, tipos, formatos y restricciones de validación aplicadas
por el esquema SQL y por los validadores de backend (\texttt{storage.py} y
\texttt{capture\_validation.py}).

\begin{longtable}{p{3.2cm} p{2.2cm} p{3.2cm} p{5.2cm}}
    \toprule
    \textbf{Atributo} & \textbf{Tipo SQL} & \textbf{Formato} & \textbf{Restricciones / Descripción} \\
    \midrule
    \endhead
    id & \texttt{INTEGER} & autonumérico & Clave primaria (\texttt{PRIMARY KEY AUTOINCREMENT}). \\
    client\_message\_id & \texttt{INTEGER} & timestamp ms & Identificador asignado por el frontend en el navegador (\texttt{nullable}). \\
    text & \texttt{TEXT} & string UTF-8 & Texto del mensaje señalado; restricción \texttt{CHECK (length(text) BETWEEN 1 AND 500)}. \\
    first\_reported\_at & \texttt{TEXT} & ISO 8601 UTC & Fecha y hora del primer reporte registrado (\texttt{NOT NULL}). \\
    \bottomrule
    \caption{Diccionario de datos — reported\_message (SQLite)}
    \label{tab:dat-dic-reported-msg}
\end{longtable}

\begin{longtable}{p{3.2cm} p{2.2cm} p{3.2cm} p{5.2cm}}
    \toprule
    \textbf{Atributo} & \textbf{Tipo SQL} & \textbf{Formato} & \textbf{Restricciones / Descripción} \\
    \midrule
    \endhead
    message\_id & \texttt{INTEGER} & FK entera & Clave foránea hacia \texttt{reported\_message(id)} con eliminación en cascada. \\
    position & \texttt{INTEGER} & entero $\geq 0$ & Posición secuencial de la palabra dentro de la frase; parte de la PK compuesta. \\
    word & \texttt{TEXT} & string & Etiqueta de la seña reconocida; \texttt{CHECK (length(word) BETWEEN 1 AND 32)}. \\
    \bottomrule
    \caption{Diccionario de datos — recognized\_word (SQLite)}
    \label{tab:dat-dic-rec-word}
\end{longtable}

\begin{longtable}{p{3.2cm} p{2.2cm} p{3.2cm} p{5.2cm}}
    \toprule
    \textbf{Atributo} & \textbf{Tipo SQL} & \textbf{Formato} & \textbf{Restricciones / Descripción} \\
    \midrule
    \endhead
    id & \texttt{INTEGER} & autonumérico & Clave primaria. \\
    message\_id & \texttt{INTEGER} & FK indexada & Referencia al mensaje reportado (\texttt{reported\_message(id)}). \\
    station & \texttt{TEXT} & string enumerado & Estación que reportó; \texttt{CHECK (station IN ('paciente', 'recepcion'))}. \\
    reported\_at & \texttt{TEXT} & ISO 8601 UTC & Marca temporal en la que se realizó la acción de reporte. \\
    \bottomrule
    \caption{Diccionario de datos — translation\_report (SQLite)}
    \label{tab:dat-dic-trans-report}
\end{longtable}

\begin{longtable}{p{3.2cm} p{2.2cm} p{3.2cm} p{5.2cm}}
    \toprule
    \textbf{Atributo} & \textbf{Tipo SQL} & \textbf{Formato} & \textbf{Restricciones / Descripción} \\
    \midrule
    \endhead
    id & \texttt{INTEGER} & autonumérico & Clave primaria. \\
    type & \texttt{TEXT} & enumerado & \texttt{CHECK (type IN ('prediction', 'sentence', 'report', 'consent', 'translation\_displayed', 'fin\_atencion'))}. \\
    occurred\_at & \texttt{TEXT} & ISO 8601 UTC & Marca temporal de ocurrencia; indexado junto con \texttt{type}. \\
    station & \texttt{TEXT} & enumerado & \texttt{NULL} o \texttt{'paciente'} / \texttt{'recepcion'}. \\
    consent & \texttt{TEXT} & enumerado & \texttt{NULL} o \texttt{'accepted'} / \texttt{'declined'} (cubre RF-03.1). \\
    lstm\_ms & \texttt{REAL} & milisegundos & Tiempo de inferencia de la red LSTM. \\
    confidence & \texttt{REAL} & porcentaje & Nivel de confianza retornado por el modelo. \\
    translation\_ms & \texttt{REAL} & milisegundos & Tiempo de generación de oración en lenguaje natural. \\
    tts\_ms & \texttt{REAL} & milisegundos & Tiempo de síntesis de voz Piper. \\
    total\_ms & \texttt{REAL} & milisegundos & Latencia de procesamiento completa en backend. \\
    e2e\_ms & \texttt{REAL} & milisegundos & Latencia extremo a extremo percibida en frontend ($0 \leq \mathrm{e2e} \leq 60000$). \\
    source & \texttt{TEXT} & string & Origen de la traducción (\texttt{'template'}, \texttt{'cache'}, \texttt{'llm'}). \\
    word\_count & \texttt{INTEGER} & entero $\geq 0$ & Cantidad de palabras en la frase procesada. \\
    \bottomrule
    \caption{Diccionario de datos — audit\_event (SQLite)}
    \label{tab:dat-dic-audit-event}
\end{longtable}

\begin{longtable}{p{2.8cm} p{2.2cm} p{3.2cm} p{5.6cm}}
    \toprule
    \textbf{Atributo} & \textbf{Tipo} & \textbf{Formato / forma} & \textbf{Restricciones} \\
    \midrule
    \endhead
    etiqueta & \texttt{str} & directorio & Debe cumplir la expresión regular
    \texttt{\textasciicircum{}[a-z0-9\_]\{1,32\}\$} (validado por
    \texttt{validate\_capture\_label}). \\

    timestamp\_ms & \texttt{int64} & epoch ms (nombre de archivo) &
    Generado por el backend como \texttt{int(time.time() * 1000)}; único
    dentro de la carpeta de la etiqueta; es el nombre base compartido por
    el archivo de landmarks y el de blendshapes. \\

    frames & \texttt{float32} & \texttt{(45, 75, 4)} & Mínimo 5 frames
    (\texttt{len(frames) < 5} rechaza); flujo normal de 45 frames (RF-T-03).
    Valores finitos (\texttt{np.isfinite}); 75 landmarks por frame
    con 4 componentes (\texttt{x, y, z, visibility}). \\

    blendshapes & \texttt{float32} & \texttt{(45, 52)} &
    52 coeficientes faciales de marcadores no manuales de LSA extraídos en la
    misma pasada de MediaPipe; valores finitos. \\
    \bottomrule
    \caption{Diccionario de datos — MuestraDeEntrenamiento (filesystem)}
    \label{tab:dat-dic-muestra}
\end{longtable}

\begin{longtable}{p{2.8cm} p{2.2cm} p{3.2cm} p{5.6cm}}
    \toprule
    \textbf{Atributo} & \textbf{Tipo} & \textbf{Formato / forma} & \textbf{Restricciones} \\
    \midrule
    \endhead
    indice & \texttt{int} & posición en el arreglo & Alineado con la capa de
    salida del modelo cargado (índice 0 = primera clase reconocida). \\

    etiqueta & \texttt{str} & string dentro de un array JSON &
    \texttt{models/custom\_classes.json} es un arreglo JSON de strings,
    escrito por \texttt{train\_lstm.py} antes de entrenar. \\
    \bottomrule
    \caption{Diccionario de datos — VocabularioDelModelo (filesystem)}
    \label{tab:dat-dic-vocabulario}
\end{longtable}

\begin{longtable}{p{2.8cm} p{2.2cm} p{3.2cm} p{5.6cm}}
    \toprule
    \textbf{Atributo} & \textbf{Tipo} & \textbf{Formato / forma} & \textbf{Restricciones} \\
    \midrule
    \endhead
    archivo & \texttt{str} & nombre de archivo \texttt{.pth} &
    \texttt{models/lstm.pth}: mejor checkpoint del último reentrenamiento. \\

    contenido & binario & \texttt{state\_dict} serializado &
    Escrito con \texttt{torch.save}; checkpoint de mejor exactitud de validación. \\

    frames\-Por\-Muestra & \texttt{int} & escalar &
    45 (ventana fija de reconocimiento, ver RF-01.3). \\
    \bottomrule
    \caption{Diccionario de datos — ModeloDeReconocimiento (filesystem)}
    \label{tab:dat-dic-modelo}
\end{longtable}

\begin{longtable}{p{2.8cm} p{2.2cm} p{3.2cm} p{5.6cm}}
    \toprule
    \textbf{Atributo} & \textbf{Tipo} & \textbf{Formato / forma} & \textbf{Restricciones} \\
    \midrule
    \endhead
    modelo & binario & \texttt{.onnx} &
    \texttt{voices/es\_AR-daniela-high.onnx}; modelo Piper TTS. \\

    configuración & JSON & sidecar &
    \texttt{voices/es\_AR-daniela-high.onnx.json}; metadatos fonémicos. \\
    \bottomrule
    \caption{Diccionario de datos — ArchivoDeVoz (filesystem)}
    \label{tab:dat-dic-voz}
\end{longtable}

\begin{longtable}{p{2.8cm} p{2.2cm} p{3.2cm} p{5.6cm}}
    \toprule
    \textbf{Atributo} & \textbf{Tipo} & \textbf{Formato / forma} & \textbf{Restricciones} \\
    \midrule
    \endhead
    id & \texttt{number} & identificador de mensaje &
    Campo del tipo \texttt{ChatMessage} (\texttt{web/types.ts}). \\

    role & \texttt{str} & \texttt{"PACIENTE"} \textbar\ \texttt{"RECEPCIONISTA"} &
    Indica el emisor del mensaje. \\

    text & \texttt{str} & texto libre & Contenido del mensaje. \\

    time & \texttt{str} & hora de envío & Hora local formateada. \\

    reported & \texttt{boolean} & booleano opcional & \texttt{true} si fue
    reportado por CU-10; sincronizado por \texttt{BroadcastChannel}. \\
    \bottomrule
    \caption{Diccionario de datos — Conversación / Mensaje (memoria volátil)}
    \label{tab:dat-dic-conversacion}
\end{longtable}

% ----------------------------------------------------------------------------
\subsection{Mapeo físico}
\label{sec:dat-mapeo-fisico}

La Tabla~\ref{tab:dat-mapeo-fisico} consolida el mapeo exhaustivo entre el
modelo conceptual, el esquema relacional lógico y su implementación concreta
en disco o memoria.

\small
\begin{longtable}{>{\raggedright\arraybackslash}p{4.0cm} >{\raggedright\arraybackslash}p{4.8cm} >{\raggedright\arraybackslash}p{5.0cm}}
    \toprule
    \textbf{Entidad conceptual} & \textbf{Destino físico / Implementación} & \textbf{Formato y motor} \\
    \midrule
    \endhead
    MensajeReportado &
    \texttt{data/lsa\_salud.db} $\rightarrow$ tabla \texttt{reported\_message} &
    SQLite (relacional con WAL) \\

    PalabraReconocida &
    \texttt{data/lsa\_salud.db} $\rightarrow$ tabla \texttt{recognized\_word} &
    SQLite (relacional con clave foránea) \\

    ReporteDeTraducción &
    \texttt{data/lsa\_salud.db} $\rightarrow$ tabla \texttt{translation\_report} &
    SQLite (relacional indexado) \\

    EventoDeAuditoría (incluye \texttt{consent}, RF-03.1) &
    \texttt{data/lsa\_salud.db} $\rightarrow$ tabla \texttt{audit\_event} &
    SQLite (relacional indexado por tipo y fecha) \\

    MuestraDe\allowbreak Entrenamiento (landmarks) &
    \texttt{dataset/\allowbreak custom/\allowbreak <etiqueta>/\allowbreak <timestamp\_ms>.npy} &
    NumPy binario \texttt{float32 (45, 75, 4)} \\

    MuestraDe\allowbreak Entrenamiento (blendshapes) &
    \texttt{dataset/\allowbreak custom/\allowbreak <etiqueta>/\allowbreak <timestamp\_ms>.blendshapes.npy} &
    NumPy binario \texttt{float32 (45, 52)} \\

    VocabularioDelModelo &
    \texttt{models/\allowbreak custom\_classes.json} &
    JSON plano (arreglo de strings) \\

    ModeloDe\allowbreak Reconocimiento &
    \texttt{models/\allowbreak lstm.pth} &
    PyTorch serializado (\texttt{state\_dict}) \\

    ArchivoDeVoz &
    \texttt{voices/es\_AR-daniela-high.onnx} + \texttt{.json} &
    ONNX runtime + metadata JSON \\

    Caché de traducciones y audio &
    Memoria del proceso FastAPI &
    \texttt{OrderedDict} LRU en RAM (volátil) \\

    Conversación ordinaria &
    Memoria del navegador (\texttt{useChatChannel}) &
    React State vía \texttt{BroadcastChannel} (volátil) \\
    \bottomrule
    \caption{Mapeo entre entidades lógicas y persistencia física}
    \label{tab:dat-mapeo-fisico}
\end{longtable}
